% 第9章: 実践・センサーモジュールを作る
\section{実践 — センサーモジュールを作る}

この章では、温度センサーモジュールを一から作ります。前章で実コードの読み方を確認したので、今度は自分で書いてみましょう。

%% ── Q1: 新モジュールの作り方 ──────────────────
\begin{qabox}{新しいモジュールはどう作るの？}

新しいモジュールを作る手順は、いつも同じ7ステップです。

\begin{enumerate}
  \item Config構造体を定義する
  \item Data構造体を定義する
  \item クラスを宣言する（.hファイル）
  \item メソッドを実装する（.cppファイル）
  \item SystemDataにDataを追加する
  \item ProjectConfig.hにConfigインスタンスを追加する
  \item main.cppのモジュール配列に追加する
\end{enumerate}

順番に見ていきましょう。analogRead()で温度を読む簡易センサーを題材にします。

\subsection{ステップ1: Config構造体}

起動時に決まる設定をまとめます。

\begin{lstlisting}
// TempSensorModule.h 内で定義
struct TempSensorConfig {
    uint8_t analogPin;         // アナログ入力ピン
    float conversionFactor;    // ADC値→温度の変換係数
};
\end{lstlisting}

\subsection{ステップ2: Data構造体}

実行中に変化するデータをまとめます。デフォルト値を必ず設定します。

\begin{lstlisting}
// TempSensorModule.h 内で定義
struct TempSensorData {
    float temperature = -999.0f;  // 温度（未取得時は異常値）
    bool isValid = false;         // データが有効か
};
\end{lstlisting}

\begin{notebox}[なぜデフォルト値を設定するの？]
init()前やセンサー故障時でも、SystemDataの値は参照される可能性があります。デフォルト値がないと、不定値（ゴミデータ）が入ったまま使われてしまいます。\texttt{-999.0f}のような明らかにおかしい値を入れておけば、「まだ読めていない」ことが一目で分かります。
\end{notebox}

\subsection{ステップ3: クラス宣言（.hファイル）}

\begin{lstlisting}
// TempSensorModule.h
#pragma once
#include <Arduino.h>
#include "IModule.h"
#include "ModuleTimer.h"

struct TempSensorConfig { /* (上で定義済み) */ };
struct TempSensorData   { /* (上で定義済み) */ };

struct SystemData;  // 前方宣言

class TempSensorModule : public IModule {
private:
    TempSensorConfig _config;
    ModuleTimer _readTimer;

public:
    TempSensorModule(const TempSensorConfig& config);
    bool init() override;
    void updateInput(SystemData& data) override;
};
\end{lstlisting}

ポイント:
\begin{itemize}
  \item \texttt{\#pragma once} — 二重インクルード防止
  \item \texttt{struct SystemData;} — 前方宣言（SystemData.hはインクルードしない）
  \item \texttt{: public IModule} — IModuleを継承
  \item \texttt{override} — 親クラスの仮想関数を上書きすることを明示
\end{itemize}

\subsection{ステップ4: メソッド実装（.cppファイル）}

\begin{lstlisting}
// TempSensorModule.cpp
#include "TempSensorModule.h"
#include "SystemData.h"  // ここでフル定義をインクルード

TempSensorModule::TempSensorModule(const TempSensorConfig& config)
    : _config(config) {}

bool TempSensorModule::init() {
    // アナログピンはpinMode不要（Arduino標準動作）
    _readTimer.setTime();
    Serial.println("[TempSensor] init OK");
    return true;
}

void TempSensorModule::updateInput(SystemData& data) {
    // 1秒ごとに温度を読み取る
    if (_readTimer.getNowTime() < 1000) return;
    _readTimer.setTime();

    int raw = analogRead(_config.analogPin);
    data.temp.temperature = raw * _config.conversionFactor;
    data.temp.isValid = true;
}
\end{lstlisting}

\subsection{ステップ5: SystemDataに追加}

\begin{lstlisting}
// SystemData.h
#include "LedModule.h"
#include "TempSensorModule.h"

struct SystemData {
    LedData led;
    TempSensorData temp;   // 追加！
};
\end{lstlisting}

\subsection{ステップ6: ProjectConfig.hに追加}

\begin{lstlisting}
// ProjectConfig.h
#include "TempSensorModule.h"

const TempSensorConfig TEMP_CONFIG = {
    .analogPin = 34,
    .conversionFactor = 3.3f / 4095.0f * 100.0f,
};
\end{lstlisting}

\subsection{ステップ7: main.cppに追加}

\begin{lstlisting}
// main.cpp
TempSensorModule tempSensor(TEMP_CONFIG);

IModule* inputModules[] = {
    &tempSensor,   // 入力モジュール配列に追加
};
\end{lstlisting}

以上で完成です。新しいモジュールを追加するときも、同じ7ステップを繰り返すだけです。

\end{qabox}

%% ── Q2: ModuleTimer ──────────────────────
\begin{qabox}{ModuleTimerはどう使うの？}

組み込みプログラムでありがちなのが、\texttt{delay()}を使った待ち時間の実装です。

\begin{badexample}{delay()で待つ}
\begin{lstlisting}
void updateInput(SystemData& data) {
    int raw = analogRead(_config.analogPin);
    data.temp.temperature = raw * _config.conversionFactor;
    delay(1000);  // ここで1秒間、何もできない！
}
\end{lstlisting}
\end{badexample}

\texttt{delay(1000)}を呼ぶと、プログラム全体が1秒間停止します。その間、他のモジュールも動けません。LEDの点滅もディスプレイの更新もすべて止まります。

\begin{goodexample}{ModuleTimerで周期実行}
\begin{lstlisting}
class TempSensorModule : public IModule {
private:
    ModuleTimer _readTimer;  // タイマーをメンバ変数に

public:
    bool init() override {
        _readTimer.setTime();  // 現在時刻を記録
        return true;
    }

    void updateInput(SystemData& data) override {
        // 前回から1秒経っていなければ即リターン
        if (_readTimer.getNowTime() < 1000) return;
        _readTimer.setTime();  // 時刻をリセット

        // 1秒経ったときだけ実行
        int raw = analogRead(_config.analogPin);
        data.temp.temperature = raw * _config.conversionFactor;
        data.temp.isValid = true;
    }
};
\end{lstlisting}
\end{goodexample}

ModuleTimerは \texttt{millis()} をラップした軽量なタイマーです。

\begin{itemize}
  \item \texttt{setTime()} — 現在時刻を記録する
  \item \texttt{getNowTime()} — 前回のsetTime()からの経過ミリ秒を返す
\end{itemize}

delay()と違い、プログラムをブロックしません。経過時間が足りなければ即座にreturnし、他のモジュールに処理を譲ります。

\begin{cautionbox}[delay()は原則禁止]
3フェーズモデルでは、各モジュールの\texttt{updateInput()}/\texttt{updateOutput()}はできるだけ速く返る必要があります。\texttt{delay()}を使うと他のすべてのモジュールが影響を受けるため、ModuleTimerを使ったノンブロッキング方式を徹底してください。
\end{cautionbox}

\end{qabox}
